% =====================================================================
\chapter{Hướng dẫn cài đặt và sử dụng chương trình}
\label{app:huongdan}
% =====================================================================

\section{Cài đặt môi trường}

Chương trình chạy trên Windows với Python 3.12. Từ thư mục \code{SOURCES}, tạo môi trường ảo và cài
các thư viện cần thiết:

\begin{lstlisting}[style=plain]
py -3.12 -m venv .venv
.\.venv\Scripts\python.exe -m pip install -r requirements.txt
\end{lstlisting}

Bộ dữ liệu Linnaeus 5 được đặt trong thư mục \path{SOURCES/data/linnaeus5/}, gồm hai thư mục con
\code{train} và \code{test}, mỗi thư mục chứa năm lớp \code{berry}, \code{bird}, \code{dog},
\code{flower}, \code{other}. Cặp ảnh mẫu \code{cover.jpg} và \code{secret.png} nằm trong
\path{SOURCES/data/sample/}.

\section{Chạy kiểm thử và các kịch bản thực nghiệm}

\begin{lstlisting}[style=plain]
# 347 ca kiem thu tu dong
.\.venv\Scripts\pytest.exe -q
# Kich ban A (Hinh 4-5): ket qua trong results/single/
.\.venv\Scripts\python.exe -m scripts.run_single
# Kich ban B (Bang 2): ket qua trong results/*.csv
.\.venv\Scripts\python.exe -m scripts.run_benchmark --secret-size 156 156 --cache-chaotic-mask
\end{lstlisting}

Các tùy chọn dòng lệnh chính:
\begin{itemize}
  \item \code{--secret-size H W}: thay đổi kích thước ảnh bí mật (mặc định giữ $246\times300$);
  \item \code{--cache-chaotic-mask}: tính sẵn mặt nạ Chaotic, chỉ đo thời gian phép XOR;
  \item \code{--vary-key-per-image}: dùng khóa khác nhau cho mỗi ảnh trong benchmark;
  \item \code{--limit N}: chỉ chạy $N$ ảnh đầu tiên (chạy thử nhanh, chỉ \code{run\_benchmark});
  \item \code{--cover} (chỉ \code{run\_single}), \code{--secret}, \code{--config}: chỉ định ảnh phủ,
        ảnh bí mật, tệp cấu hình.
\end{itemize}

\section{Ứng dụng desktop}

\begin{lstlisting}[style=plain]
# chay giao dien
.\.venv\Scripts\python.exe app.py
# kiem tra khoi dong (tu dong dong cua so)
.\.venv\Scripts\python.exe app.py --smoke
# dong goi thanh dist\CT204_CryptoStego\
.\.venv\Scripts\pyinstaller.exe packaging\ct204_app.spec
\end{lstlisting}

Bản đóng gói chỉ cần giải nén thư mục \code{CT204\_CryptoStego} và chạy
\code{CT204\_CryptoStego.exe}, không cần cài Python. Ảnh stego luôn được lưu ở định dạng PNG; hệ thống
từ chối lưu JPEG vì nén mất dữ liệu sẽ phá hủy các bit LSB.

\section{Tái tạo hình ảnh và số liệu trong báo cáo}

Các hình và số liệu bổ sung trong Chương~\ref{chap:vandung} được sinh bởi các script trong thư mục
\path{Reports/My_reports/Latex/tools}. Chúng chạy bằng môi trường ảo của dự án và không sửa đổi mã
nguồn lõi:

\begin{lstlisting}[style=plain]
cd SOURCES
.\.venv\Scripts\python.exe ..\Reports\My_reports\Latex\tools\make_figures.py
.\.venv\Scripts\python.exe ..\Reports\My_reports\Latex\tools\make_theory_figures.py
.\.venv\Scripts\python.exe ..\Reports\My_reports\Latex\tools\run_experiments.py
\end{lstlisting}
